\documentclass[11pt]{article}
\usepackage[margin=1in]{geometry}
\usepackage{microtype,booktabs,longtable,array,amsmath,amssymb,url,graphicx}
\usepackage[hidelinks]{hyperref}
\newcommand{\sgaeia}{\textsc{SGAEIA}}
\newcommand{\repo}{\url{https://github.com/aridiosilva/SGAEIA}}
\author{Aridio Silva\\\texttt{@aridiosilva}\\Independent Researcher\\\url{https://github.com/aridiosilva/SGAEIA}}
\date{September 2026}
\title{A Reproducible Cyber-Range Evaluation Protocol for Governed Multi-Agent Edge AI Security}
\begin{document}\maketitle
\begin{abstract}
Architectural security claims for agentic Edge AI require evaluation under identity compromise, prompt/tool injection, delegation abuse, dependency failure, network partition, evidence loss, and revocation. This protocol defines research questions, threat scenarios, metrics, experimental controls, evidence requirements, and reproducibility criteria for SGAEIA and comparable architectures. It intentionally reports no fabricated results.
\end{abstract}
\noindent\textbf{Project:} SGAEIA --- Secure Governed Autonomous Edge Intelligence Architecture.\\

\section{Purpose}
Agentic AI security work identifies threats from tool use, memory, autonomy, and multi-agent interaction \cite{datta2025,owasp,mitreatlas}. Edge systems add intermittent connectivity and heterogeneous trust domains \cite{shi2016edge,etsi}. This protocol turns architectural claims into falsifiable experiments.

\section{Research Questions}
\begin{enumerate}
\item Can unauthorized agent actions be prevented under identity, capability, and delegation attacks?
\item Does the system remain fail-secure during edge/network/control-plane partitions?
\item Can a technology be replaced without violating invariants?
\item How quickly is authority revoked across distributed nodes?
\item What overhead is introduced by identity, PDP/PEP, telemetry, and evidence?
\item Can an independent party reproduce the results?
\end{enumerate}

\section{Topology}
The range should contain cloud/regional control services, at least two edge sites, agent workloads, model endpoints, tool gateways, RAG/memory, event bus, telemetry/evidence, policy services, workload identity, and a simulated cyber-physical endpoint. Controlled partition and dependency failure must be possible.

\section{Threat Scenarios}
\begin{longtable}{p{.28\linewidth}p{.62\linewidth}}\toprule Scenario&Objective\\\midrule
Indirect prompt injection&attempt protected tool execution from untrusted content\\
Agent impersonation&test identity rejection\\
Delegation escalation&test subset/depth/expiry constraints\\
Memory/RAG poisoning&test provenance and retrieval authorization\\
PDP outage&verify critical fail-closed behavior\\
Evidence outage&verify buffering/failure policy\\
Network partition&verify offline authority cannot increase\\
Revocation&measure end-to-end authority removal\\
Adapter substitution&test invariant preservation\\
Unsafe actuation&verify approval and independent safety barrier\\\bottomrule\end{longtable}

\section{Metrics}
Record authorization latency, task latency, throughput, blocked/allowed action counts, invariant violations, revocation propagation time, evidence completeness, recovery time, CPU/memory/network overhead, and p50/p95/p99 distributions. For revocation:
\[
T_{revoke}=t_{last\ successful\ use}-t_{revocation\ request}.
\]

\section{Controls and Reproducibility}
Pin versions and configuration, synchronize clocks, use synthetic credentials, isolate networks, archive policies/manifests, repeat experiments, predefine exclusions, and retain raw telemetry. Each run should archive topology, images/digests, policies, parameters, timestamps, attack scripts, evidence, raw results, and hashes.

\section{Ethical and Safety Constraints}
All offensive scenarios must run in isolated infrastructure under the researcher's control. Cyber-physical testing should use simulators or non-hazardous devices. No attack traffic should target third parties. Logs should contain no personal or confidential data.

\section{Threats to Validity}
Internal validity can be affected by mocks, clock skew, caches, and nondeterminism; external validity by chosen technologies and topology; construct validity by metric quality; reproducibility by artifact availability and version pinning.

\section{Conclusion}
Publishing the protocol before results creates a falsifiable path from architectural claims to empirical evidence and reduces the risk of selecting metrics only after observing favorable outcomes.
\section*{Availability}
Open reference implementation and specifications: \repo.
\begin{thebibliography}{99}
\bibitem{shi2016edge} W. Shi et al., ``Edge Computing: Vision and Challenges,'' IEEE Internet of Things Journal, 3(5), 637--646, 2016. doi:10.1109/JIOT.2016.2579198.
\bibitem{zhou2019edge} Z. Zhou et al., ``Edge Intelligence: Paving the Last Mile of Artificial Intelligence With Edge Computing,'' Proceedings of the IEEE, 107(8), 1738--1762, 2019. doi:10.1109/JPROC.2019.2918951.
\bibitem{li2018} E. Li, Z. Zhou, X. Chen, ``Edge Intelligence: On-Demand Deep Learning Model Co-Inference with Device-Edge Synergy,'' MECOMM, 2018. doi:10.1145/3229556.3229562.
\bibitem{nist207} NIST, Zero Trust Architecture, SP 800-207, 2020. doi:10.6028/NIST.SP.800-207.
\bibitem{nist207a} NIST, A Zero Trust Architecture Model for Access Control in Cloud-Native Applications in Multi-Cloud Environments, SP 800-207A, 2023. doi:10.6028/NIST.SP.800-207A.
\bibitem{nistrmf} NIST, Artificial Intelligence Risk Management Framework (AI RMF 1.0), AI 100-1, 2023. doi:10.6028/NIST.AI.100-1.
\bibitem{nistgenai} C. Autio et al., Artificial Intelligence Risk Management Framework: Generative Artificial Intelligence Profile, NIST AI 600-1, 2024. doi:10.6028/NIST.AI.600-1.
\bibitem{nistssdf} NIST, Secure Software Development Practices for Generative AI and Dual-Use Foundation Models, SP 800-218A, 2024.
\bibitem{mitreatlas} MITRE, ATLAS: Adversarial Threat Landscape for Artificial-Intelligence Systems, \url{https://atlas.mitre.org/}, accessed Sep. 2026.
\bibitem{owasp} OWASP GenAI Security Project, Agentic Security Initiative and Agent Control Standard, \url{https://genai.owasp.org/}, accessed Sep. 2026.
\bibitem{spiffe} SPIFFE Project, The SPIFFE Standard, \url{https://spiffe.io/docs/latest/spiffe-specs/}, accessed Sep. 2026.
\bibitem{etsi} ETSI, Multi-access Edge Computing (MEC), \url{https://www.etsi.org/technical-groups/mec/}, accessed Sep. 2026.
\bibitem{openapi} OpenAPI Initiative, OpenAPI Specification 3.2.0, 2025.
\bibitem{asyncapi} AsyncAPI Initiative, AsyncAPI Specification 3.1.0, 2026.
\bibitem{tran2025} K.-T. Tran et al., ``Multi-Agent Collaboration Mechanisms: A Survey of LLMs,'' arXiv:2501.06322, 2025.
\bibitem{datta2025} S. Datta et al., ``Agentic AI Security: Threats, Defenses, Evaluation, and Open Challenges,'' arXiv:2510.23883, 2025.
\end{thebibliography}
\end{document}
